\documentclass[10pt,a4paper]{amsart}
\usepackage[margin=19mm]{geometry}
\usepackage{amsmath,amssymb,mathtools}
\usepackage[scr=rsfs,cal=euler]{mathalfa}
\usepackage{xfrac}
\usepackage[unicode,hidelinks,bookmarks=false]{hyperref}
\newcommand{\ie}{i.e.\ }
\newcommand{\cf}{cf.\ }
\newcommand{\resp}{respectively\ }
\newcommand{\Subsection}{\subsection}
\providecommand{\nobreakdash}{\nobreak\hbox{-}}
\setlength{\emergencystretch}{2em}
\multlinegap0pt
\usepackage{amssymb}
\usepackage{tikz}
\usepackage[shortlabels]{enumitem}
\usepackage{xcolor}
\newtheorem{theorem}{Theorem}
\newtheorem{claim}{Claim}
\newcommand{\Sym}{\mathop{\mathrm{Sym}}}
\title{On the approximation of permutons}
\author{Bal\'azs Maga}
\date{Source: arXiv:2605.02298v1 (CC BY 4.0)}
\begin{document}
\maketitle

\noindent\emph{Source and licence.} On the approximation of permutons. Version arXiv:2605.02298v1 (CC BY 4.0), distributed by arXiv under the Creative Commons Attribution 4.0 International licence, http://creativecommons.org/licenses/by/4.0/, as recorded in the arXiv licence metadata for this exact version and shown on the abstract page https://arxiv.org/abs/2605.02298v1. Full licence notice: \texttt{licenses/arxiv-2605.02298-CC-BY-4.0.txt}.\par\smallskip

\noindent\textit{Editorial scope.} Counted unit: the article's theorem thm:only\_bijections\_are\_well\_approximable (source0.tex lines 201--203). The article's complete original proof of it is delivered with it (source0.tex lines 420--469, one \texttt{proof} environment of 50 lines, which the article places away from the statement). No mathematics is written, repaired or re-derived: the statement and the proof are the article's own lines, in the article's own notation, and no retained line is substituted or re-flowed.  No further source-local context is needed: every cross-reference of every retained line resolves inside the two delivered spans.  Nothing was omitted from any retained span, so the package carries no omission record.  No retained line contains a citation, so the wrapper carries no bibliography and no bibliography entry.  No retained line contains a figure environment, an image-inclusion command or drawing code, so no image is delivered and the package declares no asset dependency.  Editorial formatting only, in addition: an AMS \texttt{amsart} preamble that restates the article's own declarations and macros used by the retained lines, namely the environments \texttt{theorem}, \texttt{claim} on separate counters, together with the standard packages those lines need.  The retained environments print in this document as Theorem 1, Claim 1 in order of appearance, whereas the article prints the same items with the numbering of its own full document.  The complete native source of the article as distributed in the arXiv e-print for this exact version is bundled unchanged as \texttt{sources/199709/source0.tex}.
\begin{theorem}\label{thm:only_bijections_are_well_approximable}
    Assume that $D_n(\mu)=o(1/n)$ along a subsequence for a permuton $\mu$. Then there exists a measure-preserving bijection $f:[0,1]\to[0, 1]$ defined almost everywhere such that $\mu=\mu_f$, i.e. $\mu$ is the uniform measure on the graph of $f$.
\end{theorem}

\begin{proof}[Proof of Theorem \ref{thm:only_bijections_are_well_approximable}]
    Fix a subsequence $n_1, n_2, \dots$ in accordance with the assumption of the theorem, that is $D_{n_k}(\mu)<\alpha_{k}/n_k$ with $\alpha_{k}\to 0$.
    By discarding certain entries of the sequence, we may assume that $\alpha_{k}$ is decreasing and $\sum_{k=1}^{\infty} \alpha_{k} < \infty$. We may also assume that $\frac{n_{k+1}}{n_k}>2^k$. To ease notation, set $\mu_{k} = \widehat{\pi_{k}}$ with $\pi_{k}\in \Sym(n_k)$ realizing $D_{n_k}$.

     Recall that a base square of $\pi_{k}$ is a grid square of size $n_k^{-1}\times n_k^{-1}$ contained by the support of $\widehat{\pi_k}$. For $x\in [0, 1]\setminus \mathbb{Q}$, let $Q_k(x)=I_k(x)\times J_k(x)$ be the unique base square of $\pi_{k}$ for which $x\in I_k(x)$. (For rational $x$, this would not be unique.) We say that the sequence $Q_k(x)$ is \emph{eventually nested} if $Q_{k+1}(x)\subseteq Q_k(x)$ for sufficiently large $k$.

    \begin{claim}\label{claim:stab}
        For almost every $x\in [0, 1]$, $Q_k(x)$ is eventually nested.
    \end{claim}

    \begin{proof}[Proof of Claim \ref{claim:stab}]
        Consider a base square $Q_k = I_k \times J_k$ of $\pi_k$. Define the rectangles $R_k^1, R_k^2$ to be non-overlapping with $Q_k$ such that $$I_k\times [0, 1] = R_k^1 \cup Q_k\cup R_k^2.$$
        Plugging in these rectangles into the definition of $d_\square$ verifies that for $i=1, 2$
        \begin{equation}\label{eq:claim_stab}
        \mu(R_k^i)\leq \frac{\alpha_k}{n_k}.
        \end{equation}
        Now consider base squares of $\pi_{k+1}$ over $I_k$, the $x$-projections of these are $1/n_{k+1}$ grid intervals. As $n_k$ may not be a divisor of $n_{k+1}$, these grid intervals may not be compatible with the $1/n_k$ grid interval $I_k$, however, $I_k$ fully contains at least $\frac{n_{k+1}}{n_k}-2$ of them. Among the corresponding base squares there can be at most 2 which intersects $Q_k$ nontrivially, all others are either contained in $R_k^1$, $R_k^2$, or $Q_k$. Observe that $R_k^i$ ($i=1, 2$) can contain at most $2\alpha_k \frac{n_{k+1}}{n_k}$ of them. Indeed, otherwise in $R_k^i$, the difference of $\mu$ and $\mu_{k+1}$ would be at least $\frac{\alpha_k}{n_k}$ due to \eqref{eq:claim_stab}, contradicting $d_\square(\mu_{k+1}, \mu)\leq \frac{\alpha_{k+1}}{n_{k+1}}$. To sum up, we conclude the following for $Q_{k+1}(x)$ if $x\in I_k$ (and hence $Q_k = Q_k(x)$:
        \begin{itemize}
            \item on a set of measure at most $\frac{2}{n_{k+1}}$, we have no control over $Q_{k+1}(x)$, this measure corresponds to the overhanging parts of the $1/n_{k+1}$ grid intervals.
            \item on a set of measure at most $\frac{2}{n_{k+1}}$, $Q_{k+1}(x)$ partially intersects $Q_k(x)$, hence it is not contained in it.
            \item on a set of measure at most $\frac{4\alpha_k}{n_k}$, $Q_{k+1}$ is contained in either $R_k^1$ or $R_k^2$, hence it is disjoint from $Q_k$.
            \item in the remainder of $I_k$, $Q_{k+1}(x)\subseteq Q_k(x)$
        \end{itemize}
        Hence in $I_k$, $Q_{k+1}(x)\subseteq Q_k(x)$ apart from measure at most $\frac{4}{n_{k+1}} + \frac{4\alpha_k}{n_k}$.
        Multiplying by $n_k$ to take into account all possible choices of $I_k$, $Q_{k+1}(x)\subseteq Q_k(x)$ fails in $[0,1]$ at a set of measure at most
        $$\frac{4n_k}{n_{k+1}} + 4\alpha_k.$$
        By the assumptions on $(\alpha_k)$ and $(n_k)$, the resulting series is summable, thus by the Borel--Cantelli lemma, we have that $Q_{k+1}(x)\subseteq Q_k(x)$ holds for almost every $x$ for sufficiently large $k$. This proves the claim.
    \end{proof}

    For almost every irrational $x$, whose existence is guaranteed by the claim, denote by $k(x)$ the minimal index from where $Q_k(x)$ is nested. Since the intervals \(J_j(x)\) are nested and have lengths \(1/n_j\to0\), their intersection consists of a single point; define \(f(x)\) to be this point.
     
    Observe that the function \(f\) is injective outside \(f^{-1}(\mathbb Q)\). Indeed, for $x\neq y$ and sufficiently large $K>k(x), k(y)$, we have $J_K(x)\neq J_K(y)$, thus $f(x)=f(y)$ might happen only if these values coincide with the shared endpoint of $J_K(x)$ and $J_K(y)$.
    
    After defining \(f\) arbitrarily on the remaining null set, let \(\mu_f\) be the pushforward of the Lebesgue measure under \(x\mapsto(x,f(x))\). We claim that it is the weak limit of the sequence $\mu_{k}$. This would conclude the proof: as the limit is unique, we obtain $\mu_f = \mu$, which directly implies that $f$ is a measure-preserving function, and by omitting the null-set $f^{-1}(\mathbb{Q})$ from its domain, we get a measure-preserving bijection due to our observation about $f$ being almost injective.

    By the Portmanteau theorem, it suffices to show that for any open set $G\subseteq [0, 1]^2$ we have
    \begin{equation}\label{eq:portmanteau}
    \liminf_{k\to\infty} \mu_k(G)\geq \mu_f(G).
    \end{equation}
    Fix $G$, and let $H=\{x:\ (x, f(x))\in G\}$. If $\lambda(H)=0$, then $\mu_f(G)=0$ and we have nothing to prove. Otherwise for any $\varepsilon>0$ for every large enough integer $K$ we can find $H'\subseteq H$ with the properties 
    \begin{itemize}
    \item $\lambda(H')>\lambda(H)-\varepsilon$,
    \item for any $x\in H'$, the square with sides $\frac{2}{n_K}$ centered at $(x, f(x))$ is contained in $G$,
    \item for any $x\in H'$, $k(x)\leq K$.
    \end{itemize}
    (By taking $K\to \infty$, both the second and third properties hold for each $x\in H$, which can be replaced by a finite $K$ with arbitrarily small loss of measure.)

    Now it is clear that $Q_K(x)\subseteq G$ for $x\in H'$ and $H'\subseteq \bigcup_{x\in H'} I_K(x)$. This yields that there are at least $\lambda(H')n_K$ distinct intervals in this union, or equivalently, at least $\lambda(H')n_K$ distinct base squares of $\pi_K$ in $G$. As $\mu_K$ puts measure $1/n_K$ to each such square, $\mu_K(G)\geq \lambda(H)-\varepsilon$. This proves \eqref{eq:portmanteau}, hence concludes the proof of the theorem.
    
\end{proof}

\end{document}
